\chapter{Los Diez Arquetipos}

Los arquetipos del sistema \textbf{ZINA} proporcionan la lente humana a través de la cual los jugadores experimentan NeoBilbao. No son meras clases mecánicas: son posiciones en la cadena de montaje, trincheras en el asfalto y fracturas morales en un mundo despiadado.

A continuación se presentan los diez arquetipos oficiales: cinco nacidos en la lluvia ácida de \textbf{2084} y cinco forjados bajo el Velo y el metal azul de \textbf{2102}.

\section{Arquetipos de 2084: La Fractura Temprana}

\begin{FichaArquetipo}{Arquetipo 1: El Sabueso (2084)}
\textbf{«En la Margen Izquierda no hay crímenes sin resolver, solo crímenes que no compensa cobrar».}

Exagente de seguridad despedido, cobrador de deudas callejero o detective privado a sueldo de los comités vecinales. Conoces los callejones donde las farolas parpadean en ámbar y sabes qué precio tiene cada confidencia en las tabernas de Barakaldo.

\textbf{Emociones sugeridas:} Rencor $+2$, Arraigo $+1$, Fatiga $0$, Esperanza $-1$.

\textbf{Movimientos de Arquetipo:}
\begin{itemize}
    \item \textbf{Olfato de Asfalto:} Cuando investigas la escena de un altercado o desaparición usando \textit{Arraigo}, obtienes 1 reserva de pista adicional incluso con un resultado de $7\text{--}9$.
    \item \textbf{El Mal Asunto:} Cuando interrogas a un testigo reacio apelando al miedo o su pasado con \textit{Rencor}, puedes forzarle a responder con la verdad sin que alerte a la policía corporativa, a cambio de marcar 1 casilla de Zama.
    \item \textbf{Revólver de Fogueo y Plomo:} Tu arma corta casera nunca se encasquilla en espacios reducidos. Si disparas a quemarropa, sumas $+1$ al daño cinético infligido.
    \item \textbf{Red de Soplones:} Una vez por sesión, puedes declarar que conoces a un borracho, camarero o basurero local que te debe un favor clave sobre el caso en curso.
\end{itemize}

\textbf{Equipo Inicial:} Gabardina encerada y raída, revólver pesado calibre .38 con tambor modificado (\textit{Daño 2, Corto, Estruendosa}), grabadora analógica de microcinta, petaca de orujo casero y libreta con direcciones tachadas.
\end{FichaArquetipo}

\vspace{1em}

\begin{FichaArquetipo}{Arquetipo 2: La Biohacker (2084)}
\textbf{«La carne enferma y los chips se queman; mi trabajo es que ambos sigan funcionando hasta el viernes».}

Cirujana clandestina, artesana de ciberware de desguace y experta en terminales CRT. Sabes perforar un cráneo para implantar un sintonizador ilegal en la trastienda de una carnicería tanto como puentear un relé de Orekak-Corp.

\textbf{Emociones sugeridas:} Esperanza $+2$, Fatiga $+1$, Rencor $0$, Arraigo $-1$.

\textbf{Movimientos de Arquetipo:}
\begin{itemize}
    \item \textbf{Bisturí y Soldador:} Cuando operas de urgencia a un compañero para curar una Herida Grave, tira con \textit{Esperanza}. Con $10+$, la herida baja a Leve de inmediato sin consumir suministros quirúrgicos de farmacia.
    \item \textbf{Overclocking Neumático:} Puedes sobrecargar temporalmente un implante mecánico o una herramienta pesada. Otorga ventaja inmediata a la siguiente acción, pero con $7\text{--}9$ el dispositivo chisporrotea y queda inoperativo hasta repararse.
    \item \textbf{Empalme Clandestino:} Cuando te conectas físicamente a un terminal CRT de la empresa mediante cable coaxial, no dejas rastro digital en los servidores centrales si superas la tirada con $10+$.
    \item \textbf{Farmacopea Callejera:} Posees una reserva de viales de analgésicos sintetizados y estimulantes que pueden ignorar penalizadores de heridas durante 1 escena.
\end{itemize}

\textbf{Equipo Inicial:} Maletín de instrumental quirúrgico desinfectado con aguardiente, soldador de estaño portátil, monitor CRT de muñeca con cables de pinza de cocodrilo, pistola neumática de sutura y 3 viales de coagulante.
\end{FichaArquetipo}

\vspace{1em}

\begin{FichaArquetipo}{Arquetipo 3: El Celador (2084)}
\textbf{«Los de arriba tienen tanquetas blindadas; nosotros nos tenemos los unos a los otros».}

Obrero de choque, portero de asamblea vecinal o vigilante nocturno de los silos de carbón. Eres el escudo humano de tu comunidad; cuando las furgonetas de los Secretas aparcan en la esquina, eres tú quien se planta en mitad de la calzada con una barra de encofrar.

\textbf{Emociones sugeridas:} Fatiga $+2$, Rencor $+1$, Arraigo $0$, Esperanza $-1$.

\textbf{Movimientos de Arquetipo:}
\begin{itemize}
    \item \textbf{Muro de Hormigón:} Cuando te interpones físicamente para proteger a un aliado de un ataque o explosión, tiras con \textit{Fatiga}. Con $10+$, absorbes todo el daño reduciéndolo en 1 nivel sin caer derribado.
    \item \textbf{Grito de Barricada:} Cuando lideras un grupo de vecinos o trabajadores desarmados frente a fuerzas de choque, gastas 1 punto de Auzolan para que la multitud no se disperse ante gases lacrimógenos ni disparos de advertencia.
    \item \textbf{Puño de Remache:} Tus golpes cuerpo a cuerpo con herramientas pesadas o puños de acero cuentan con la etiqueta \etiqueta{Contundente, Demoledora}.
    \item \textbf{Tozudez de Hierro:} Una vez por sesión, puedes negarte a caer inconsciente tras sufrir una Herida Crítica hasta haber puesto a salvo a un compañero.
\end{itemize}

\textbf{Equipo Inicial:} Barra de encofrar reforzada con empuñadura de cuero (\textit{Daño 2, Melé, Contundente}), chaleco de kevlar caducado de segunda mano (\textit{Armadura 1}), linterna industrial de carburo y silbato de asamblea de latón.
\end{FichaArquetipo}

\vspace{1em}

\begin{FichaArquetipo}{Arquetipo 4: La Burócrata (2084)}
\textbf{«Una bala puede matar a un hombre; un formulario con el sello adecuado desahucia a una manzana entera».}

Antigua archivista municipal o empleada de administración en los escalafones bajos de Orekak-Corp. Conoces los códigos de barras, los conductos de expedientes, los nombres de los jueces corruptos y los sellos de lacre digital que abren compuertas blindadas.

\textbf{Emociones sugeridas:} Esperanza $+1$, Arraigo $+2$, Rencor $0$, Fatiga $-1$.

\textbf{Movimientos de Arquetipo:}
\begin{itemize}
    \item \textbf{Papeles en Regla:} Puedes falsificar pases de transporte, salvoconductos de toque de queda o permisos de obra en minutos. Con $10+$, ningún guardia de control cuestionará la autenticidad del documento.
    \item \textbf{Fuga en el Archivo:} Cuando buscas información confidencial en expedientes impresos o microfichas, utilizas \textit{Arraigo}. Sabes exactamente qué funcionario aceptará un soborno o qué cajón no tiene cerradura.
    \item \textbf{Sello Corporativo:} Tienes acceso a una firma electrónica o sello físico en desuso de Orekak-Corp que puede retrasar una redada o clausura policial durante 24 horas.
    \item \textbf{Oído en los Despachos:} Cuando te mezclas entre mandos intermedios y secretarios corporativos, puedes sonsacar órdenes inminentes sin levantar sospechas.
\end{itemize}

\textbf{Equipo Inicial:} Portafolios de cuero rígido con plantillas de sellado oficial, microcámara analógica de bolsillo, pistola de bolsillo Derringer (\textit{Daño 1, Discreta}), identificación de empleada de bajo nivel y tarjeta perforada maestra.
\end{FichaArquetipo}

\vspace{1em}

\begin{FichaArquetipo}{Arquetipo 5: El Familiar (2084)}
\textbf{«Dijeron que se lo llevó la Lotería, pero mientras yo respire, nadie se da por muerto».}

Madre, hermano, cónyuge o hijo de alguien que fue subido a los transportes de Orekak-Corp y nunca volvió a dar señales de vida. Tu motor no es el dinero ni el poder; es una llama obstinada de amor y rabia que la policía y los auditores no comprenden ni pueden apagar.

\textbf{Emociones sugeridas:} Esperanza $+2$, Arraigo $+1$, Rencor $0$, Fatiga $-1$.

\textbf{Movimientos de Arquetipo:}
\begin{itemize}
    \item \textbf{Retrato en el Bolsillo:} Cuando contemplas el recuerdo de tu ser querido antes de tomar una decisión desesperada, puedes eliminar 1 casilla de Zama o sumar $+1$ a tu siguiente tirada de \textit{Afrontar la Presión}.
    \item \textbf{Llamada de Sangre:} Si un compañero de cuadrilla está en peligro de muerte, puedes moverte hacia él a través de fuego cruzado o escombros ignorando penalizadores de movimiento.
    \item \textbf{Empatía Desarmante:} Cuando apelas a la humanidad o la familia de un guardia, trabajador corporativo o sicario, tiras con \textit{Esperanza}. Con $10+$, el enemigo duda, baja el arma o te permite pasar discretamente.
    \item \textbf{La Voz de los Sin Nombre:} Puedes convocar en cuestión de minutos a otros familiares de desaparecidos para montar una concentración de protesta improvisada que sirva de distracción ruidosa.
\end{itemize}

\textbf{Equipo Inicial:} Fotografía analógica de tu familiar en un relicario de estaño, llave de tubo o cuchillo de cocina afilado (\textit{Daño 1}), linterna de dinamo, megáfono de mano a pilas y libreta con nombres de desaparecidos.
\end{FichaArquetipo}

\section{Arquetipos de 2102: El Siglo Blindado}

\begin{figure}[htbp]
  \centering
  \includegraphics[width=\textwidth]{foundry_workers_platform.jpg}
  \caption{Turno de relevo en los hornos de Burdin-Urdin: la vida bajo la disciplina del Velo.}
\end{figure}

\begin{FichaArquetipo}{Arquetipo 6: La Fundidora (2102)}
\textbf{«El acero azul no respeta los títulos; si no sabes manejar la cuchara de colada, te disuelve igual que a una escoria».}

Piloto de exoesqueleto de carga pesada Basajaun-MK3 en los hornos de colada continua de Sestao. Tu cuerpo está curtido por el calor insoportable y el peso del metal. Para ti, doblar vigas de aleación y romper barricadas blindadas es solo otra parte de la jornada laboral.

\textbf{Emociones sugeridas:} Fatiga $+2$, Rencor $+1$, Arraigo $0$, Esperanza $-1$.

\textbf{Movimientos de Arquetipo:}
\begin{itemize}
    \item \textbf{Piel de Amianto:} Tu entrenamiento y tu mono ignífugo te otorgan inmunidad al fuego ordinario y reducen en 2 puntos cualquier daño por calor, plasma o salpicaduras de metal líquido.
    \item \textbf{Llave de Trinquete:} En combate cuerpo a cuerpo o maniobras con maquinaria pesada, puedes sustituir \textit{Rencor} por \textit{Fatiga} al tirar los dados.
    \item \textbf{Carga Basajaun:} Cuando arremetes contra una puerta blindada, un muro de hormigón o un vehículo con tu exoesqueleto, tiras con \textit{Fatiga}. Con $10+$, abres una brecha transitable de un solo impacto.
    \item \textbf{Brazos de Remache:} Puedes levantar, sostener o doblar pesos industriales colosales durante unos instantes para permitir que tus compañeros escapen de un derrumbe.
\end{itemize}

\textbf{Equipo Inicial:} Exoesqueleto de carga Basajaun desgastado (+1 Fuerza física, \etiqueta{Aparatoso}), soplete térmico de corte de plasma (\textit{Daño 3, Penetrante, Fuego}), máscara de soldador con filtro UV integrado y guantes de kevlar aluminizado.
\end{FichaArquetipo}

\vspace{1em}

\begin{FichaArquetipo}{Arquetipo 7: El Técnico de Mantenimiento (2102)}
\textbf{«Los oligarcas construyen sus torres sobre planos que dibujamos nosotros; conocemos cada conducto y cada válvula de purga».}

Especialista en las cloacas mecánicas, conductos de refrigerante de Burdin-Urdin y subestaciones eléctricas del Zócalo. Te mueves por las tripas de la ciudad donde los drones Tártalo no tienen señal y los oligarcas temen mirar.

\textbf{Emociones sugeridas:} Esperanza $+2$, Fatiga $+1$, Arraigo $0$, Rencor $-1$.

\textbf{Movimientos de Arquetipo:}
\begin{itemize}
    \item \textbf{Atajo por los Conductos:} Cuando tu cuadrilla huye de una patrulla corporativa o drones de seguridad, puedes guiarles a través de tuberías y falsos techos tirando con \textit{Fatiga}. Con $10+$, desaparecéis sin dejar rastro térmico.
    \item \textbf{Cortocircuito Táctico:} Puedes manipular una caja de distribución eléctrica o una consola del Velo en combate para cegar sensores, apagar luces o electrificar charcos de agua a tu alrededor.
    \item \textbf{Ojo de Tártalo:} Posees un dron de reconocimiento Tártalo modificado y reprogramado que puede volar en sigilo para explorar habitaciones o hackear terminales a distancia.
    \item \textbf{Parche de Emergencia:} Puedes reparar temporalmente cualquier mecanismo averiado, ascensor o compuerta neumática con alambre y chatarra en segundos.
\end{itemize}

\textbf{Equipo Inicial:} Multímetro de frecuencias del Velo, dron de bolsillo con cámara óptica, llave de grifa de tres bocas (\textit{Daño 1, Contundente}), rollo de cable de fibra óptica y arnés de escalada industrial.
\end{FichaArquetipo}

\vspace{1em}

\begin{FichaArquetipo}{Arquetipo 8: La Delegada Sindical (2102)}
\textbf{«Si una sola máquina se detiene, la Cúpula se oscurece. Esa es la única verdad que entienden».}

Portavoz elegida de las asambleas secretas de taller, organizadora de paros de brazos caídos y defensora comunitaria. Conoces el Estatuto Laboral de Orekak-Corp mejor que sus propios abogados y sabes cómo inflamar el coraje de mil obreros exhaustos con una sola frase.

\textbf{Emociones sugeridas:} Arraigo $+2$, Rencor $+1$, Esperanza $0$, Fatiga $-1$.

\textbf{Movimientos de Arquetipo:}
\begin{itemize}
    \item \textbf{Voz de Megáfono:} Cuando arengas a una cuadrilla desmoralizada o a un grupo de huelguistas, tiras con \textit{Arraigo}. Con $10+$, el grupo recupera 2 puntos de Auzolan o disipa 1 punto de Zama colectiva.
    \item \textbf{Piquete de Asfalto:} Puedes ordenar el cierre coordinado de un paso de peatones, montacargas o compuerta de sector mediante operarios leales, impidiendo el paso de refuerzos de seguridad durante 30 minutos.
    \item \textbf{Caja de Resistencia:} Tienes acceso a suministros clandestinos (raciones alimentarias, munición reciclada, tarjetas de crédito opacas) custodiados por el sindicato.
    \item \textbf{El Peso del Convenio:} Cuando amenazas a un capataz o auditor de bajo rango con paralizar un sector clave si no cede a una demanda inmediata, tiras con \textit{Rencor}. Con $10+$, el mando corporativo transige para evitar sanciones de sus superiores.
\end{itemize}

\textbf{Equipo Inicial:} Megáfono neumático de bolsillo, copia digital encriptada de los convenios laborales restringidos, porra telescópica de acero (\textit{Daño 2, Discreta}), brazalete rojo sindical y un sintonizador del Velo con bloqueador de geolocalización.
\end{FichaArquetipo}

\vspace{1em}

\begin{FichaArquetipo}{Arquetipo 9: El Rebotado (2102)}
\textbf{«En el Yermo no hay sirenas de fábrica ni ojos del Velo; solo polvo cáustico y libertad para morir como tú elijas».}

Exiliado, desertor de turnos forzosos o contrabandista que cruza los muros perimetrales hacia el Yermo contaminado. Conoces los senderos entre las ruinas de las antiguas autovías, comercias con chatarra militar prohibida y tus pulmones han desarrollado resistencia a los vapores químicos.

\textbf{Emociones sugeridas:} Rencor $+2$, Fatiga $+1$, Esperanza $0$, Arraigo $-1$.

\textbf{Movimientos de Arquetipo:}
\begin{itemize}
    \item \textbf{Pulmones de Plomo:} Eres completamente inmune a las toxinas ambientales y gases industriales estándar de baja concentración sin necesidad de máscara respiratoria.
    \item \textbf{Carroñero de Ruinas:} Cuando buscas equipo, piezas de repuesto o munición entre los restos de un sector clausurado o en el Yermo, tiras con \textit{Fatiga}. Con $10+$, encuentras exactamente la pieza exótica o el cargador que la cuadrilla necesitaba.
    \item \textbf{Tiro Sucio:} Cuando atacas por sorpresa desde las sombras o emboscada, puedes añadir la etiqueta \etiqueta{Brutal} a tu ataque y sumar $+2$ al daño infligido.
    \item \textbf{Sendero Oculto:} Conoces brechas en los muros de contención de la ciudad por donde cruzar mercancías o fugitivos burlando los radares de la Cúpula.
\end{itemize}

\textbf{Equipo Inicial:} Rifle de cerrojo modificado con mira telescópica casera (\textit{Daño 3, Largo alcance, Recarga manual}), poncho impermeable con camuflaje de hollín, máscara antigás militar de filtro regenerativo, kit de ganzúas y brújula magnética analógica.
\end{FichaArquetipo}

\vspace{1em}

\begin{FichaArquetipo}{Arquetipo 10: La Médica de Cuadrilla (2102)}
\textbf{«El Burdin-Urdin solidifica la sangre si te descuidas; mi bisturí decide quién vuelve a fichar mañana».}

Facultativa de dispensario obrero, veterana de quirófanos de campaña en los barracones o desertora de los laboratorios sanitarios de la Cúpula. Tu lucha es contra la muerte silenciosa que traen los vapores de la fundición, el agotamiento extremo y la temida Plaga Gris.

\textbf{Emociones sugeridas:} Esperanza $+2$, Arraigo $+1$, Fatiga $0$, Rencor $-1$.

\textbf{Movimientos de Arquetipo:}
\begin{itemize}
    \item \textbf{Triaje de Urgencia:} En medio del combate o un desastre industrial, puedes estabilizar a un aliado al borde de la muerte con una tirada de \textit{Esperanza}. Con $10+$, el personaje no muere y puede continuar en pie con 1 casilla de salud funcional.
    \item \textbf{Inyección de Choque:} Administras un cóctel de adrenalina y quelantes de Burdin-Urdin que permite a un aliado ignorar todo dolor y actuar como si tuviera salud plena durante una escena completa (tras la cual sufre 1 casilla de Zama).
    \item \textbf{Purga de Velo:} Puedes realizar una punción lumbar rápida para desconectar temporalmente el nodo neural de un obrero de la red del Velo, impidiendo que los capataces rastreen sus constantes biométricas durante 6 horas.
    \item \textbf{Diagnóstico Frío:} Puedes determinar con una simple mirada la causa exacta de una herida, veneno o infección tecnológica, descubriendo si proviene de armamento corporativo o sabotaje.
\end{itemize}

\textbf{Equipo Inicial:} Inyector neumático de alta presión con 5 dosis de quelante azul, maletín quirúrgico de campo con desinfectante, bisturí térmico (\textit{Daño 1, Penetrante}), escáner biométrico portátil y bata médica reforzada con lona plomada.
\end{FichaArquetipo}
